% Candidate section. This file is intended for inclusion in the root paper.
% IDs F-SM-* identify results proved in this run, not claims of historical novelty.
\section{The two supplied routes: exact scope and mathematical bridges}
\label{sec:source-masullo}

\subsection{What is adopted, what has been checked, and what is not exported}

We use Alp\"oge's five declarations in \texttt{ALPOGE\_Main.lean} as the research
inputs specified in the task.  The check performed here is an inspection of the
declarations and of \texttt{ALPOGE\_Statements.lean}, not a fresh Lean kernel
verification: the import \texttt{Supplements.Kernels.MainTheoremsAll} is absent
from the supplied ZIP, as is the remainder of the formal development.  Neither
\texttt{lean} nor \texttt{lake} is installed in the working environment.  The
supplied scan is a report about another source tree; no claim is made that the
175,187 files were rescanned here.  This distinction does not change the adopted
mathematical premises, but it determines the correct check-event label.

The declaration names and their exact targets are as follows.
\begin{enumerate}
\item \texttt{riemannHypothesis : RiemannHypothesis := rhOfAll}; the target is
Mathlib's named proposition.
\item \texttt{riemannZeta\_eq\_zero\_re : RiemannZetaZerosOnLine := rhOfAll};
the supplied definition excludes the negative even integers and $s=1$.
\item \texttt{LFunction\_ne\_zero\_of\_half\_lt\_re :
DirichletLNonvanishingRightOfHalf := allChi}; this quantifies over every positive
modulus, every complex Dirichlet character, and $\Re s>1/2$, with the exception
$(\chi,s)=(1,1)$.
\item \texttt{LFunction\_eq\_zero\_re : DirichletLZerosOnLine := claimAll.2.1};
this concerns zeros in the open strip $0<\Re s<1$.
\item \texttt{LFunction\_eq\_zero\_re\_of\_isPrimitive :
PrimitiveDirichletLZerosOnLine := critAll}; this is the primitive-character
specialization, again in the open strip.
\end{enumerate}
The pinned Mathlib revision is
\texttt{0df444a360eaa60ab8c11dca51a86af692955474}, in
\texttt{ALPOGE\_lakefile.toml}.  None of these five declarations states a
weight-monodromy theorem, a theorem on all Hecke characters of all number
fields, a nonabelian Artin holomorphy theorem, or a theorem specifying an
integral algebraic $K$-group of $\mathbb Z$.

Two corrections to the prompt's reading of the four-page note are material.
The theta-lift chain has length of order $\sqrt X$, both in the paragraph on
page~1 and in step~8 on page~4.  The radical is visible in the PDF.  The claim
about $K$-groups in step~17, footnote~6, is explicitly tentative and gives no
degrees or quantifiers.  It cannot be turned into a specified $K$-group
vanishing theorem merely by citing the note.  The relevant export, if one
exists in the full development, must first be identified.

\subsection{Masullo's different target and its exact remaining gate}

Masullo, \emph{The Music of the Primes}, page~6, specifies all number fields
and all unitary Hecke characters.  This is a larger target than the five
Alp\"oge declarations above.  The analytic route can be located precisely:
\begin{enumerate}
\item Section~1.2.1, pages~13--15: a divisor exact sequence of nuclear
Fr\'echet spaces, then Hardy--Sobolev Hilbert completions whose spectral
quotients retain zeros in a prescribed vertical strip.
\item Sections~1.2.2--1.2.3, pages~15--16, and Chapter~11: orbital bases
and a decomposition of a lifted geometric kernel into local terms.
\item Theorem~10.4.2.8, pages~463--471: the spectral trace after the
relevant truncation/projection operations.
\item Theorem~11.2.2.1, pages~490--491: orbital nuclearity on the
zero line for Sobolev degree $s\geq 2[F:\mathbb Q]$.
\item Theorem~11.3.2.1, pages~495--496: orbital trace calculations;
Theorem~12.4.1.1, pages~514--519: the $S$-local trace calculation.
\item Theorem~12.5.1.1, pages~519--520: the proposed global trace
identity with only the zeros inside the chosen strip on its spectral side.
\item Sections~1.4.2--1.4.5, pages~28--30: the proposed deduction of
zero location by holomorphic families of test functions.
\end{enumerate}
The end of Chapter~13 is not a usable theorem statement: on page~522,
Theorem~13.3.1.2 consists of a heading with instructions to insert the
statement and prove the main results.  Theorem~13.3.1.1 on the same page is
marked for correction.  More substantially, Section~12.3 on page~514, devoted
to the global-to-$S$-local deconstruction, consists only of its heading and an
instruction to include the section.  The proof of Theorem~12.5.1.1 on
page~520 uses the resulting decomposition and exchange of diagonal integration
with the local sum.  These are concrete missing links in the supplied written
argument.  The following proposition supplies the zero-location deduction
once the appropriate trace identity has been established; it does not supply
the missing kernel identity.

\paragraph{Finding F-SM-01 (a shrinking-strip trace criterion).}
Let $\mathcal Z$ be a locally finite multiset in
$\{\lambda\in\mathbb C:|\Re\lambda|\leq B\}$, with positive integral
multiplicities $m(\lambda)$, and suppose
\[
 \sum_{\lambda\in\mathcal Z}
       \frac{m(\lambda)}{1+|\Im\lambda|^2}<\infty.
\]
For $f\in C_c^\infty(0,\infty)$ put
$\widehat f(\lambda)=\int_0^\infty f(t)e^{\lambda t}\,dt$.
Assume that for some sequence $a_j\downarrow0$,
\begin{equation}\label{eq:sm-strip}
 \sum_{\lambda\in\mathcal Z}m(\lambda)\widehat f(\lambda)
  =\sum_{\substack{\lambda\in\mathcal Z\\|\Re\lambda|<a_j}}
           m(\lambda)\widehat f(\lambda)
 \qquad(f\in C_c^\infty(0,\infty)).
\end{equation}
Then every $\lambda\in\mathcal Z$ is imaginary.  Conversely, if every
$\lambda$ is imaginary, \eqref{eq:sm-strip} holds for every $a_j>0$.

\begin{proof}
Fix $a=a_j$ and let $\mathcal E$ be the submultiset with
$|\Re\lambda|\geq a$.  The hypothesis gives
$\sum_{\mathcal E}m(\lambda)\widehat f(\lambda)=0$ for every such test
function.  Fix a real number $w>B$ and a complex number $z$ with
$\Re z>B$.  On $t\geq0$ define
\[
 h_{z,w}(t)=e^{-zt}-e^{-wt}.
\]
Choose smooth functions $\eta,\zeta$ on $[0,\infty)$ with $0\leq\eta,\zeta
\leq1$, $\eta=0$ on $[0,1]$, $\eta=1$ on $[2,\infty)$, $\zeta=1$ on
$[0,1]$, and $\zeta=0$ on $[2,\infty)$.  Set
$f_{\epsilon,R}(t)=\eta(t/\epsilon)\zeta(t/R)h_{z,w}(t)$, where
$0<\epsilon\leq1$ and $R\geq1$.  These are admissible compactly supported
smooth test functions on $(0,\infty)$.

For $\beta\in[-B,B]$, the $L^1$ norms of
$g_{\epsilon,R,\beta}=e^{\beta t}f_{\epsilon,R}(t)$ and its second
derivative are bounded independently of $\epsilon,R,\beta$.  Here is the
estimate at the only potentially singular endpoint.  Since $h_{z,w}(0)=0$,
$h_{z,w}(t)=O(t)$, $h'_{z,w}(t)=O(1)$ and $h''_{z,w}(t)=O(1)$ for
$0\leq t\leq2$.  On $[\epsilon,2\epsilon]$ the terms containing
$\eta''(t/\epsilon)/\epsilon^2$ therefore have $L^1$ norm
$O(\epsilon^{-2}\epsilon\epsilon)=O(1)$, and those containing
$\eta'(t/\epsilon)/\epsilon$ also have $L^1$ norm $O(1)$.
The remaining terms near zero are bounded.  At infinity all terms are
bounded by a polynomial in $t$ times $e^{-ct}$, with
$c=\min(\Re z,w)-B>0$; the derivatives of $\zeta(t/R)$ have size
$O(R^{-1})$ and $O(R^{-2})$.  This proves the uniform bounds.
Twice integrating by parts in the Fourier variable, together with the
$L^1$ bound when $|\gamma|\leq1$, now gives
\[
 |\widehat f_{\epsilon,R}(\beta+i\gamma)|
 \leq\frac{C}{1+\gamma^2}
 \qquad(\beta\in[-B,B]).
\]
Dominated convergence in $t$, followed by dominated convergence in the
absolutely summable zero series, gives
\[
 0=\sum_{\lambda\in\mathcal E}m(\lambda)
       \left(\frac1{z-\lambda}-\frac1{w-\lambda}\right)
       \qquad(\Re z>B).
\]
The series on the right is normally convergent on every compact subset of
$\mathbb C\setminus\mathcal E$: outside a finite set each summand is bounded
by $C_Km(\lambda)/(1+|\Im\lambda|^2)$.  It therefore defines a
meromorphic function on $\mathbb C$.  At any $\lambda_0\in\mathcal E$
its residue is the strictly positive integer $m(\lambda_0)$; local finiteness
ensures that the other terms are holomorphic near $\lambda_0$.
The identity theorem, applied to the connected complement of the discrete
set $\mathcal E$, makes the meromorphic function identically zero.  A zero
meromorphic function has no pole, so $\mathcal E$ is empty.  This holds for
every $a_j$, proving $\Re\lambda=0$ for every $\lambda$.  The converse
is immediate because no terms are omitted.
\end{proof}

\paragraph{Finding F-SM-02 (the precise theorem-level comparison).}
For one fixed unitary Hecke character, suppose the established full explicit formula
and a proposed shrinking-strip trace formula have exactly the same
geometric distribution on test functions supported at idelic norm $>1$.
After taking its character isotype and setting $t=\log|g|$, their equality
is \eqref{eq:sm-strip} for $\lambda=\rho-1/2$.  The usual zero-counting
bound $N(T)=O(T\log(T+2))$ implies the summability required in
F-SM-01, so the trace formula implies GRH for this character.  Conversely,
GRH implies this numerical trace identity by the full explicit formula.
For the Dirichlet characters covered by Alp\"oge's adopted theorem, the
numerical identity is therefore available without assuming Masullo's
operator-theoretic construction.

\begin{proof}
The trace distributions subtract to exactly \eqref{eq:sm-strip}.  By a
dyadic decomposition,
\[
 \sum_{2^k<|\Im\lambda|\leq2^{k+1}}
 \frac{m(\lambda)}{1+|\Im\lambda|^2}
 \ll \frac{2^k(k+1)}{2^{2k}},
\]
whose sum is finite; local finiteness handles bounded ordinates.
Apply F-SM-01.  Conversely, all zeros lie in every open strip of positive
width centered on the critical line, so restricting the spectral sum removes
no terms.  This proves the equality of distributions.  No assertion about
the existence of a particular kernel, basis, or geometric realization is
used or obtained in this converse.
\end{proof}

The full idelic explicit formula is recorded in Masullo's Theorem~B.4.4.7,
page~623, whereas the restricted one is Theorem~12.5.1.1,
pages~519--520.  To apply F-SM-02 to the manuscript one must first align the
Haar-measure and multiplicity conventions in those statements, and prove
the missing deconstruction.  F-SM-01 is independent of those normalization
choices.  It also avoids passing separately to divergent archimedean
integrals: the difference of exponentials vanishes at the identity.

\subsection{A larger Artin bridge, conditional on Masullo's larger target}

\paragraph{Finding F-SM-03 (off-line Artin holomorphy from Hecke GRH).}
Assume GRH for every finite-order Hecke character of every number field.
Let $K$ be a number field and $V$ a finite-dimensional complex
representation of $\operatorname{Gal}(\overline K/K)$ with finite image.
Then its Artin $L$-function is holomorphic and nonzero on
\[
 U=\{s:0<\Re s<1,\quad\Re s\ne1/2\}.
\]
The same assertion for $V(n)$ holds on $U-n$ under the convention
$L(V(n),s)=L(V,s+n)$.

\begin{proof}
Choose a finite Galois extension $E/K$ through which $V$ factors and put
$G=\operatorname{Gal}(E/K)$.  Brauer induction gives
\[
 [V]=\sum_{j=1}^r n_j\operatorname{Ind}_{H_j}^{G}\psi_j,
 \qquad n_j\in\mathbb Z,\qquad\dim\psi_j=1.
\]
The $\psi_j$ have finite order.  With $K_j=E^{H_j}$, Artin reciprocity
and inductivity of Artin $L$-functions give the meromorphic identity
\[
 L_K(V,s)=\prod_{j=1}^r L_{K_j}(\psi_j,s)^{n_j}.
\]
Each Hecke factor is holomorphic in $0<\Re s<1$, including the trivial
character's Dedekind zeta function, whose possible pole is at $1$.
The assumed GRH says that no such factor vanishes on $U$.  An integral
power, positive, zero, or negative, of a nowhere-zero holomorphic function
is again nowhere-zero and holomorphic.  Their finite product has the same
property.  Substituting $s+n$ proves the twisted assertion.
\end{proof}

The exact representation-theoretic input is Serre (1977),
\emph{Linear Representations of Finite Groups}, \S10.5, Theorem~20,
page~78.  The $L$-function factorization is equation~(2) in Cogdell,
\emph{On Artin $L$-functions}, \S2.1, page~10; the compatibility of the
ramified local factors is explained in \S2.2, page~12.  This is a conditional
bridge to all finite-image Artin motives, substantially beyond the
abelian-over-$\mathbb Q$ scope of the supplied Alp\"oge theorem.
It does not prove holomorphy on the critical line: factors in the denominator
can vanish there.  Nor does the proof imply general automorphy, the Artin
conjecture, or the existence of a motivic realization functor.

\subsection{Two explicit corrections to the geometric analogies}

\paragraph{Finding F-SM-04 (the norm does not recover the unit strata).}
The assertion on Masullo page~493 that a saturated all-local-units stratum
in $\mathbb A_F^\times$ is the inverse image of the corresponding unit
stratum of $\mathbb A_{\mathbb Q}^\times$ under the norm map is false.
Here ``saturated'' has its explicit meaning from page~492: a set of places
contains all places over a rational prime whenever it contains one.
The displayed stratum on page~493 is a subset of the ideles themselves,
defined by $g_v\in\mathcal O_{F_v}^{\times}$ at each listed place;
it is not a subset enlarged by multiplication by $F^{\times}$.
The correct condition at a rational prime $p$ for the norm to be a unit is
\[
 \sum_{v\mid p} f(v/p)v_v(g_v)=0,
\]
where $v_v$ assigns value $1$ to a uniformizer and $f(v/p)$ is the residue
degree.  The stronger all-local-units condition is $v_v(g_v)=0$ for
every $v\mid p$.

\begin{proof}
For a finite extension $F_v/\mathbb Q_p$, a uniformizer has norm of
$p$-adic valuation $f(v/p)$, and the norm of a unit is a unit; by factoring
an arbitrary element into a uniformizer power times a unit,
$v_p(N_{F_v/\mathbb Q_p}(g_v))=f(v/p)v_v(g_v)$.  Multiplication of the
local norm factors proves the displayed formula.
Take $F=\mathbb Q(i)$ and $p=5$, which splits into two degree-one places
$v_1,v_2$.  Thus $F_{v_j}\simeq\mathbb Q_5$.  Let the idele $g$ have
$g_{v_1}=5$, $g_{v_2}=5^{-1}$ and every other component $1$.
Its norm is the identity idele because the two $5$-adic components multiply
to $1$.  Neither of these components is a local unit, so $g$ does not lie
in the asserted unit stratum.  This is the required counterexample.
The all-local-units stratum remains well-defined directly; the error in
identifying it with a norm inverse image does not by itself disprove
the later trace identity.
\end{proof}

\paragraph{Finding F-SM-05 (a maximal curve need not have gonality $g+1$).}
The equality for maximal curves asserted on Masullo page~535 is false.
Over $\mathbb F_9$, the smooth plane curve
\[
 H:\quad Y^3Z+YZ^3=X^4
\]
is maximal of genus $3$ and has gonality at most $3$, hence strictly
less than $g+1=4$.

\begin{proof}
In characteristic $3$, the three partial derivatives of
$Y^3Z+YZ^3-X^4$ are $-X^3,Z^3,Y^3$, so they have no common projective
zero.  The curve is smooth.  Over an algebraic closure, any two positive-degree
plane-curve components intersect by B\'ezout's theorem, creating a singular
point; a repeated component would also contradict smoothness.  Thus this
smooth plane quartic is geometrically integral.  Its genus is
$(4-1)(4-2)/2=3$ (the plane-curve genus formula, Stacks Project,
Lemma~53.9.3).  For each $x\in\mathbb F_9$, $x^4$
is the norm of $x$ and lies in $\mathbb F_3$; the trace
$y\mapsto y^3+y$ from $\mathbb F_9$ onto $\mathbb F_3$ has three
elements in each fibre.  There are therefore $9\cdot3=27$ affine
points.  At $Z=0$ the equation forces $X=0$, leaving the unique point
$[0:1:0]$.  Thus $\#H(\mathbb F_9)=28=9+1+2\cdot3\cdot3$,
which is the maximal value.  The nonconstant rational function $x=X/Z$
has function-field degree at most $3$, because $y=Y/Z$ satisfies
$y^3+y=x^4$ over $\mathbb F_9(x)$.  It extends to a morphism to
$\mathbb P^1$ on the smooth projective curve; consequently the gonality
is at most $3$.
\end{proof}

The correction strengthens rather than defeats the point being made in
Masullo \S14.2, pages~534--536: the proposed number-field multiplicity
mechanism cannot simply be transported to function fields.  For this
curve the Frobenius eigenvalue multiplicity in $H^1$ is $2g=6$, whereas
the gonality is at most $3$: the finite-field Weil theorem gives modulus
$3$ for each of its six eigenvalues, and their sum is
$9+1-28=-18$, so equality in the triangle inequality makes all six
equal to $-3$.

\subsection{The categorical and positivity information still required}

The stack diagrams are especially relevant to Grothendieck's programme,
but their stated status matters.  Masullo's equations~(8.4.0.6) and
(8.4.0.7), page~390, and (8.4.0.8)--(8.4.0.9), page~391, relate
Fr\'echet-space divisor sequences to periodic-orbit quotient stacks and
boundary divisors.  Remark~8.4.0.1 on page~391 explicitly says that the
vertical arrows are not morphisms.  They express an interpretation, not
a constructed sheaf, realization functor, or equivalence of categories.
Consequently these diagrams do not supply an arithmetic cohomology theory
of $\operatorname{Spec}\mathbb Z$, a six-functor formalism, or a comparison
with the \emph{D\'erivateurs}.  A literal categorical bridge would require
those arrows to be constructed and their compatibility with the trace,
tensor products, Tate twists, and local terms proved.

\paragraph{Finding F-SM-06 (spectral location alone does not supply
semisimplicity or a positive unitary realization).}
There exists a complex finite-dimensional operator with purely imaginary
spectrum that is not skew-adjoint for any positive definite Hermitian
inner product.

\begin{proof}
On $\mathbb C^2$ take $A=i\gamma I+N$, where $\gamma\in\mathbb R$,
$Ne_2=e_1$ and $Ne_1=0$.  Its spectrum is the singleton $\{i\gamma\}$.
If $A$ were skew-adjoint for a positive definite Hermitian form, then
$e^{tA}$ would preserve its norm for every real $t$: differentiating
$\langle e^{tA}v,e^{tA}v\rangle$ gives zero.  But
$e^{tA}e_2=e^{i\gamma t}(e_2+te_1)$, whose norm is at least
$|t|\|e_1\|-\|e_2\|$ by the triangle inequality and therefore tends
to infinity.  This contradicts norm preservation.
\end{proof}

This refutes a general inference, not RH and not the possibility of a
particular Hilbert--P\'olya realization.  Masullo's spectral discussion
expressly permits Jordan blocks (\S1.2.1, page~14), so zero location
must be separated from a theorem killing nilpotent parts or proving
zero simplicity.

\paragraph{Finding F-SM-07 (a bounded biorthogonal lift is not automatically
a Schauder basis).}
The proof step in Masullo's Lemma~10.4.2.5, page~462, that obtains
the Schauder basis property from the displayed boundedness and
biorthogonality alone is invalid.

\begin{proof}
Let $E=\ell^2(\{0,1,2,\ldots\})$, $E_0=\mathbb Ce_0$, $F=E$,
and let $x_n$ be the image of $e_n$ in $E/E_0$ for $n\geq1$.
These $x_n$ form the usual orthonormal basis.  Set $u_n=e_0+e_n$
and let $\phi_n$ be the $n$th coordinate functional.  Then
$\phi_n(u_m)=\delta_{nm}$ and $\|u_n\|=\sqrt2$.
The collar conditions hold with $r=1$, $r'=11/10$ and $\epsilon=1/2$.
Yet
\[
 \frac1N\sum_{n=1}^{N}u_n
   =e_0+\frac1N\sum_{n=1}^{N}e_n\longrightarrow e_0
\]
in $E$.  Thus $e_0$ belongs to the closed span of the $u_n$, while
every $\phi_n(e_0)=0$.  If the $u_n$ were a Schauder basis of that
closed span, applying the continuous biorthogonal coordinate functionals
to the expansion of $e_0$ would give every coefficient zero, forcing
$e_0=0$, a contradiction.
\end{proof}

This counterexample concerns the inference used in the written proof.
It does not refute the existence of suitable lifts in the Hilbert-space
case.  The following construction supplies that repair in the form relevant
to the Hardy--Sobolev Hilbert spaces.

\paragraph{Finding F-SM-08 (Hilbert-space basis lifting with a controlled
graph).}
Let $E$ be a Hilbert space, $E_0$ a closed subspace, $q:E\to E/E_0$ the
quotient, and $F\subset E$ a linear subspace such that $F\cap E_0$ is
dense in $E_0$.  Let $(x_n)$ be a Schauder basis of $E/E_0$ with
coefficient functionals $(\phi_n)$, or a Markushevich basis with
its total biorthogonal functionals $(\phi_n)$, and assume each
$x_n\in q(F)$.
For any positive numbers $\delta_n$ satisfying
$\sum_n\|\phi_n\|\delta_n<\infty$, there are $u_n\in F$ with
$q(u_n)=x_n$, $\|u_n-Sx_n\|<\delta_n$, where
$S:E/E_0\to E_0^\perp$ is the isometric orthogonal section, such that
$(u_n)$ is respectively a Schauder or Markushevich basis of its closed span.
That span maps
isomorphically to $E/E_0$ under $q$.

\begin{proof}
Choose $v_n\in F$ with $q(v_n)=x_n$.  Then $Sx_n-v_n\in E_0$.
Density of $F\cap E_0$ provides $w_n\in F\cap E_0$ with
$\|v_n+w_n-Sx_n\|<\delta_n$; put $u_n=v_n+w_n$ and
$d_n=u_n-Sx_n\in E_0$.  The series
\[
 Kx=\sum_n\phi_n(x)d_n
\]
converges absolutely for every $x\in E/E_0$ and defines a bounded
operator to $E_0$, with norm at most
$\sum_n\|\phi_n\|\delta_n$.  Set $T=S+K$.
Then $qT=\operatorname{id}$, so $\|Tx\|\geq\|x\|$ and $T$
is a bounded linear isomorphism onto a closed subspace of $E$.
The inverse is the restriction of $q$.  Since $Tx_n=u_n$, in the
Schauder case the image under $T$ of the basis expansion of $x$ is a
convergent basis expansion in $(u_n)$.  In the Markushevich case,
density of the span and totality of the biorthogonal functionals
$\phi_n\circ q$ on the image are transported by the isomorphism $T$.
Their closed span is $T(E/E_0)$ in both cases, proving every assertion.
\end{proof}

Thus one lift-related gap has an explicit repair.  This does not repair
all of Theorem~10.4.2.8 or the missing Section~12.3: the estimates for
the resulting integral kernels and the interchange of diagonal
integration with deconstruction are separate requirements.

\subsection{Comparison of methods and what each can actually transfer}

\paragraph{Finding F-SM-09 (the exact power-sum bridge to purity).}
Let $\alpha_1,\ldots,\alpha_r$ be distinct nonzero complex numbers,
let $c_1,\ldots,c_r$ be nonzero complex coefficients, and let $R>0$.
Suppose there are constants $C>0$ and $A\geq0$ such that
\[
 \left|\sum_{j=1}^r c_j\alpha_j^n\right|\leq C n^A R^n
 \qquad(n\geq1).
\]
Then $|\alpha_j|\leq R$ for every $j$.  If additionally the set of
$\alpha_j$ is invariant under $\alpha\mapsto R^2/\alpha$, then
$|\alpha_j|=R$ for every $j$.

\begin{proof}
The power-series
\[
 F(z)=\sum_{n\geq1}\left(\sum_{j=1}^r c_j\alpha_j^n\right)z^n
\]
converges absolutely on $|z|<R^{-1}$ by the assumed bound.  Near $0$ it
is the rational function
\[
 \sum_{j=1}^r\frac{c_j\alpha_jz}{1-\alpha_jz}.
\]
For any $j$, the latter has a genuine pole at $z=1/\alpha_j$:
the $j$th summand has residue $-c_j/\alpha_j\ne0$ and the other
summands are holomorphic there because the $\alpha_j$ are distinct.
If $|\alpha_j|>R$, this pole would lie in the disk on which the power
series is holomorphic, contradicting equality by analytic continuation.
Thus $|\alpha_j|\leq R$.  Under the stated reciprocal symmetry,
$R^2/\alpha_j$ is another listed number, so
$R^2/|\alpha_j|\leq R$, which gives the reverse inequality.
\end{proof}

Taking $R=q^{w/2}$ identifies the precise missing inputs for a purity
argument of the form suggested by Alp\"oge's footnote~3: bounds over
\emph{all} finite extensions, a nonzero net coefficient at each eigenvalue,
and the appropriate duality.  A bound for a single finite field is
insufficient for this deduction.  For the trace of an honest finite-dimensional
space the coefficients are the positive algebraic multiplicities, so
they cannot cancel.  For alternating cohomological traces this must
be proved separately.

\paragraph{Finding F-SM-10 (why alternating traces alone do not give purity).}
For any $R>0$ there is a graded vector space with Frobenius whose
alternating traces satisfy every bound in F-SM-09, while a Frobenius
eigenvalue has modulus greater than $R$.

\begin{proof}
Take $V^0=V^1=\mathbb C$, all other graded pieces zero, and let
Frobenius act by the same scalar $a>R$ on both nonzero pieces.
Then $\sum_i(-1)^i\operatorname{tr}(F^n\mid V^i)=a^n-a^n=0$
for every $n\geq1$, while both pieces have eigenvalue $a$.
\end{proof}

This is the exact point of a no-cancellation or positivity input;
it supplies spectral information that a virtual determinant can hide.
Even F-SM-09 determines eigenvalue moduli, not the nilpotent part of
Frobenius, as F-SM-06 illustrates.

Alp\"oge's step~5 (\emph{Remark}, page~3 and footnote~3) seeks a
noncancelling sign in cohomological power sums; his steps~7--10 and~12--16
(page~4) obtain arithmetic multiplicity and norm estimates from unitary
theta lifts, periods, and integral structures.  Masullo's primary gain is
different: a coarser completion retains only part of the zero divisor,
while the geometric trace is asserted to remain unchanged.  F-SM-01
shows precisely why preservation of that trace forces the omitted part
to be empty.  The first route needs a positivity/divisibility estimate;
the second needs trace compatibility through completion, lifting,
localization, and a limiting process.

Neither gain automatically gives algebraic cycles, crystalline
Frobenius, a realization functor, a weight filtration on an arbitrary
motive, or reconstruction of nonabelian fundamental groups.  In the
supplied text the Iwasawa tower is an arithmetic amplification device
(Alp\"oge steps~11 and~19), and the hexagonal CM curve supplies a
specific theta-geometric setting (step~3).  Neither is identified by a
theorem with the Teichm\"uller tower or the profinite GT group.
The testable bridges are the explicit $L$-function identities and the
trace/norm compatibilities just described.  The correspondence between
the words ``tower'', ``hexagon'', or ``multiplicity'' by itself does not
constitute any such compatibility.

The Masullo route also does not eliminate the logical possibility of a
geometric or foliated realization of a different construction.  Its
Sections~1.2.7 and~1.5 compare its own analytic objects with Deninger's
expectations; the absence of a three-manifold in this construction is
not a nonexistence theorem for all other constructions.
